% Αντικατάστησε τίτλο, υπότιτλο, στόχους και περιεχόμενο με τη δική σου ύλη.
% Η BookChapter δίνει άνοιγμα κεφαλαίου χωρίς φωτογραφία.
\BookChapter{Κίνηση σε μία διάσταση}{Θέση και μετατόπιση}{%
  \item Να ορίζεις έναν προσανατολισμένο άξονα αναφοράς.
  \item Να διαβάζεις την αλγεβρική τιμή της θέσης.
  \item Να υπολογίζεις τη μετατόπιση από δύο θέσεις.
  \item Να ερμηνεύεις το πρόσημο της μετατόπισης.
}

% Προαιρετικός σύντομος τίτλος: section[κεφαλίδα/TOC]{πλήρης τίτλος}.
\section[Θέση και μετατόπιση]{Θέση και μετατόπιση}
Η θέση ενός σώματος σε έναν άξονα δίνεται από τη συντεταγμένη του.
Για κίνηση από τη θέση $x_A$ στη θέση $x_B$, η μετατόπιση είναι
\begin{equation}\label{eq:starter-displacement}
  \Delta x=x_B-x_A.
\end{equation}

\begin{ThinkBox}
Επιλέγω τη θετική φορά, σημειώνω τις δύο θέσεις και αφαιρώ την αρχική
από την τελική τιμή. Το πρόσημο δείχνει τη φορά της μετατόπισης.
\end{ThinkBox}

% Το μέγεθος των γραμμάτων και τα πάχη προέρχονται από τα κοινά styles.
% Η έμφαση είναι ανά σχήμα: εδώ τονίζεται η μετατόπιση.
\begin{figure}[htbp]
  \centering
  \begin{tikzpicture}[x=.8cm,y=.8cm]
    \draw[book/axis] (-3.2,0)--(3.5,0)
      node[book/physlabel,right] {$x\,(\mathrm{m})$};
    \foreach \x in {-2,0,2}{%
      \draw[book/tick] (\x,.08)--(\x,-.08);
      \node[book/ticklabel,below=2pt] at (\x,-.08) {$\x$};
    }
    \node[book/endpoint] at (-2,0) {};
    \node[book/endpoint] at (2,0) {};
    \node[book/physlabel,above] at (-2,0) {$A$};
    \node[book/physlabel,above] at (2,0) {$B$};
    \draw[book/mainvector] (-2,.65)--(2,.65)
      node[book/physlabel,midway,above] {$\Delta x$};
  \end{tikzpicture}
  \caption{Μετατόπιση από τη θέση Α στη θέση Β.}\label{fig:starter-displacement}
\end{figure}

Στο Σχήμα~\ref{fig:starter-displacement}, η σχέση~\eqref{eq:starter-displacement}
δίνει $\Delta x=2-(-2)=4\,\mathrm{m}$.

\begin{table}[htbp]
  \centering
  \begin{BookStripedTable}
  \begin{tabularx}{\linewidth}{@{}X>{\centering\arraybackslash}p{25mm}@{}}
    \BookTableHeaderRow
    \BookTableHead{Μέγεθος} & \BookTableHead{Τιμή}\\
    Αρχική θέση $x_A$ & $-2\,\mathrm{m}$\\
    Τελική θέση $x_B$ & $2\,\mathrm{m}$\\
    Μετατόπιση $\Delta x$ & $4\,\mathrm{m}$\\
  \end{tabularx}
  \end{BookStripedTable}
  \caption{Θέσεις και μετατόπιση του παραδείγματος.}\label{tab:starter-values}
\end{table}

\BookSection{Ερωτήσεις}
\begin{exercise}[START-001]{Σωστό ή λάθος;}
  Η μετατόπιση εξαρτάται από την αρχική και την τελική θέση.
  \TrueFalseChoices
\end{exercise}
\Answer{Σωστό.}
